% Template Latex para trabalhos academicos da PUC Minas
%
% Pontificia Universidade Catolica de Minas Gerais (www.pucminas.br)
% Instituto de Ciencas Exatas e Informatica (www.icei.pucminas.br)
%
% VERSAO: 0.2
% DATA: Maio de 2012
%
% AUTORES:
% Prof. Joao Paulo Domingos Silva
% Rodrigo Caetano de Oliveira Rocha
%
% COLABORADORES:
% Tamara Márcia Mendes
% Harison Herman Silva
%
% REVISORES:

\documentclass{abntpuc}

\usepackage[pdftex]{graphicx}
\usepackage{fancyhdr}

\begin{document}

% --- Campos que formarao: capa, folha de rosto, folha de aprovacao  ---

\curso{Bacharelado em Sistemas de informação}

\autor{Nome do Aluno}

\titulo{TÍTULO DO TRABALHO} % Caixa alta

\subtitulo{Subtítulo do Trabalho} % Caixa baixa

\cidade{Belo Horizonte}

\ano{2012}

\trabalho{Monografia} % Tipo de trabalho: Monografia, Dissertacao, Tese...

\grau{Bacharel em Ciência da Computação} % Grau do trabalho: Bacharel em..., Mestre em..., Doutor em...

\orientador{Nome do Orientador}

\avaliadorA{Nome do Avaliador 1}

\avaliadorB{Nome do Avaliador 2}

\datacompleta{DD de MM de 2012} % Dia, mes e ano

% --- CAPA ---

\capa

% --- FOLHA DE ROSTO ---

\folharosto

% --- FOLHA DE APROVACAO ---

\folhaaprovacao

% --- DEDICATORIA (elemento opcional) ---

\dedicatoria
{
	\textit{Dedicatória: Página onde o autor presta homenagem a uma ou mais pessoas.
	O layout desta página fica a critério do autor, mas o tipo e tamanho de letras são definidos pela ABNT.}
}

% --- AGRADECIMENTOS (elemento opcional) ---

\agradecimentos
{
	Agradecimentos a pessoas que contribuíram para o desenvolvimento do trabalho.
	Agradecimentos a pessoas que contribuíram para o desenvolvimento do trabalho.
}

% --- EPIGRAFE (elemento opcional) ---

\epigrafe
{
	\textit{Epígrafe: Pensamentos retirados de um livro, uma música, um poema, normalmente relacionados ao tema do trabalho.
	Deve ser elaborada conforme norma NBR 10520/2002. Apresentação de citações em documentos.
	Se desejar, a epígrafe pode ser grafada em itálico.
	Ao final do trabalho deve-se fazer a referência completa da publicação de onde a epígrafe foi retirada.}
}

% --- RESUMO ---

\resumo
{
	Apresentação concisa dos pontos relevantes do texto. Deve ressaltar o objetivo, o método, resultados e conclusões do trabalho. Deve-se utilizar o verbo na voz ativa ou terceira pessoa do singular. O resumo não deve conter citações ou indicações bibliográficas.
} % Resumo
{
	Ao final do resumo deve-se elaborar palavras-chave representativas do conteúdo do trabalho, separadas entre si por um ponto.
} % Palavras-chave

% --- ABSTRACT ---

\abstract
{
	Versão do resumo em idioma de divulgação internacional. Deve ser a tradução literal do resumo em português e apresentar palavras- chave no mesmo idioma, logo abaixo do texto, separadas entre si por um ponto.
} % Abstract
{
	O resumo em língua estrangeira também deve conter palavras-chave representativas do conteúdo do trabalho, separadas entre si por um ponto.
} % Keywords

% --- LISTA DE FIGURAS, LISTA DE TABELAS, SUMARIO ---

\listafiguras
\listatabelas
\listasiglas {
   \sigla{S1}{Sigla 1}
   \sigla{S2}{Sigla 2}
   \sigla{S3}{Sigla 3}
}
\sumario

% --- TEXTO ---

\capitulo{INTRODUÇÃO}

\iniciocapitulo
A introdução deve conter a natureza do trabalho, justificativa, objetivos, o tema proposto e outros elementos para situar o trabalho.\par

% Exemplo de tabela:

\begin{table}[!htb]
	\caption[Descrição na Lista de Tabelas]{Legenda da tabela}
	\centering
	\label{tab:tabela}
	\begin{tabular}{c|c|c}
		\hline
		Coluna 1 & Coluna 2 & Coluna 3 \\
		\hline
 		Célula 1 & Célula 2 & Célula 3 \\
		Célula 4 & Célula 5 & Célula 6 \\
		Célula 7 & Célula 8 & Célula 9 \\
		\hline
	\end{tabular}
	\\ \textbf{\footnotesize Fonte: Fonte da tabela}
\end{table}

% Exemplo de figura:

\begin{figure}[!htb]
   \caption[Descrição na Lista de Figuras]{Legenda da figura}
   \label{fig:figura1}
   \centering
   \includegraphics{LogoPUC.jpg}
   \\ \textbf{\footnotesize Fonte: Fonte da figura}
\end{figure}


\capitulo{DESENVOLVIMENTO}

\iniciocapitulo
Compreende a revisão da literatura ou estado da arte, metodologia e exposição da pesquisa.\par

A Revisão da Literatura compõe-se da evolução do tema e ideias de diferentes autores sobre o assunto. Deve conter citações textuais ou livres, com indicação dos autores conforme norma NBR 10520/2002.\par

O Estado da Arte segundo definição no dicionário Aurélio é o nível de desenvolvimento atingido (por uma ciência, uma técnica) na atualidade, ou seja, refere-se ao quadro atual de uma área, suas tendências, potencialidades e excelência no assunto.\par

A Metodologia deve apresentar o método adotado na entrevista, questionário, observação, experimentação na população pesquisada na características e quantificação.\par

A Exposição da Pesquisa é a análise dos fatos apresentados, ou seja, os dados obtidos, as estatísticas, comparações com outros estudos e outras observações.\par

\secao{Seção}

Exemplo de citação \cite{knuth68}.

Exemplo de citação \cite{knuth69}.

Exemplo de citação \cite{knuth73}.

\subsecao{Subseção}

\subsubsecao{Subsubseção}

\capitulo{CONCLUSÃO}

\iniciocapitulo
Discussão dos resultados obtidos na pesquisa, onde se verificam as observações pessoais do autor. Poderá também apresentar sugestões de novas linhas de estudo. A conclusão deve estar de acordo com os objetivos do trabalho. A conclusão não deve apresentar citações ou interpretações de outros autores.

% --- REFERENCIAS ---

% Para as referencias funcionarem, e preciso baixar e instalar o pacote ABNTEX.

\referencias{Referencias}

\end{document}
